% arara: lualatex: { branch: developer, draft: true }
% arara: lualatex: { branch: developer, draft: true }
% arara: lualatex: { branch: developer }
% arara: clean: { extensions: [ aux, log, out] }
% arara: move: { files: [ spintent-01-luamml-mathml.html ], target: 'spintent-01-mathml.html' }
\DocumentMetadata{lang=es-CL, pdfversion=2.0, pdfstandard=ua-2, tagging=on, tagging-setup={math/setup={mathml-SE}},}
\documentclass{article}
\usepackage[spanish]{babel}
\usepackage[top=2cm,bottom=2cm,left=1cm,right=1cm]{geometry}
\usepackage{unicode-math,hyperref,enumext}
\hypersetup
  {
    colorlinks = true,
    pdftitle   = {Test tagged PDF for spintent package},
  }
\usepackage{spintent}
\theoremstyle{definition}
\newtheorem{ejercicio}{Ejercicio}[section]
\begin{document}

\title{Guía Práctica: Operaciones, Notación Científica y Fracciones}
\author{Pablo González Luengo}
\date{}
\maketitle

\section{Operaciones con números decimales}

\begin{ejercicio}[Adición]
  Calcule la suma de los siguientes valores exactos, manteniendo la
  alineación correspondiente a la coma decimal:
  \[
    X = \spnum{12,45} + \spnum{3,789}
  \]
\end{ejercicio}

\begin{ejercicio}[Multiplicación]
  Determine el producto de los siguientes números decimales:
  \[
    Y = \spnum{4,2} \times \spnum{1,5}
  \]
\end{ejercicio}

\section{Conversión a notación científica}

\begin{ejercicio}[Magnitudes grandes]
  Exprese el valor numérico en notación científica utilizando el símbolo
  de multiplicación explícito (\texttt{times}):
  \[
    \spnum{45000} \implies \spnum[notation-sym=times]{4,5e4}
  \]
\end{ejercicio}

\begin{ejercicio}[Magnitudes pequeñas]
  Convierta el número a notación científica forzando la lectura del signo
  en la base:
  \[
    \spnum{0,00000000000016} \implies
    \spnum[read-sign=terse, notation-sym=cdot]{+1,6e-4}
  \]
\end{ejercicio}

\section{Conversión de decimales a fracciones}

\begin{ejercicio}[Decimal finito]
  Convierta el siguiente decimal exacto a fracción irreductible:
  \[
    \spnum{0,75} = \frac{75}{100} = \frac{3}{4}
  \]
\end{ejercicio}

\begin{ejercicio}[Decimal periódico mixto]
  Exprese el siguiente decimal periódico infinito como fracción y simplifique
  su resultado al máximo:
  \[
    \spnum{2,1;6} = \frac{216 - 21}{90} = \frac{195}{90} = \frac{13}{6}
  \]
\end{ejercicio}

\end{document}
